% 388_CMB_Grundform_En_ch.tex
% Johann Pascher, 2 Oct 2026
% The CMB temperature in the basic form: eV relation, ratio to m_e, H0 forms

\providecommand{\BM}{\textbf{[B]}}
\providecommand{\KM}{\textbf{[K]}}
\providecommand{\SM}{\textbf{[S]}}
\providecommand{\XM}{\textbf{[X]}}
\providecommand{\QM}{\textbf{[Q]}}

\chapter*{The CMB temperature in the basic form}
\addcontentsline{toc}{chapter}{The CMB temperature in the basic form}

\begin{abstract}
The Casimir-CMB connection is an identity, because the length $L_\xi$ is fixed solely through the
CMB temperature (R136). This document collects all routes meant to obtain $T_{\mathrm{CMB}}$ from
$\xi$ without this circularity and recalculates them in the basic form: natural units with
$\alpha=1$, a single anchor $m_e$. The eV relation $T_{\mathrm{CMB}}=\frac{16}{9}\xi$~eV with
second-order correction agrees to $2\times10^{-6}$, but holds only in the unit eV, whose size was
fixed in 1881 and set exactly in 2019. In the basic form, $T_{\mathrm{CMB}}/m_e$ depends on a
single number. The candidate $T_{\mathrm{CMB}}/m_e=(8\pi)^{1/4}\,\xi^{5/2}$ meets the FIRAS value
to $0.12\sigma$, together with the corpus form of $H_0$ gives $T_{\mathrm{CMB}}^4=16\,H_0\,m_e^3$,
and makes $L_\xi$ independent of the CMB. The $H_0$ form with exponent~10 fits it; the older one
with $41/4$ is off by $11\sigma$. Neither the prefactor nor the exponent is derived; what remains
open is listed at the end.
\end{abstract}

% ============================================================
\section*{Status markers}
% ============================================================
\begin{center}
\begin{tabular}{@{}lp{11cm}@{}}
\textbf{[B]} & Proven algebraically. \\[2pt]
\textbf{[K]} & Confirmed numerically (verification script, comparison with measured values). \\[2pt]
\textbf{[S]} & Setting, open or speculative. \\[2pt]
\textbf{[X]} & Refuted or not viable. \\[2pt]
\textbf{[Q]} & Convention. \\
\end{tabular}
\end{center}
Measured values: $T_{\mathrm{CMB}}=2.72548\pm0.00057$~K (FIRAS, Fixsen 2009; relative uncertainty
$0.021\,\%$), CODATA 2022. Verification script: \texttt{pruef\_388\_cmb\_grundform.py}.

% ============================================================
\section{Starting point}
% ============================================================

The length $L_\xi=(15\xi/\pi^2)^{1/4}\,\hbar c/(k_BT_{\mathrm{CMB}})=100.24\,\mu$m is chosen such
that the photon density of the CMB equals $\xi/L_\xi^4$ there. The comparison with the Casimir
energy density at the same length is therefore an identity, not a measurement (Doc.~279, R136).
An $H_0$ anchor does not change this either, since $L_\xi$ depends only on $T_{\mathrm{CMB}}$.

The $\Lambda$CDM model does not derive $T_{\mathrm{CMB}}$ either. It takes the FIRAS value as a
fixed input, and the Planck value $H_0=67.4$~km/s/Mpc is a derived quantity that presupposes it.
The question is therefore not whether FFGFT is worse off than $\Lambda$CDM, but whether it can do
more: $T_{\mathrm{CMB}}$ as a ratio from $\xi$, with the same anchor as the masses.

% ============================================================
\section{The eV relation}
% ============================================================

Since 2025 the corpus has carried the relation $T_{\mathrm{CMB}}=\frac{16}{9}\xi$~eV, later with a
second-order correction (P11):
\begin{equation}
k_BT_{\mathrm{CMB}}=\tfrac{16}{9}\,\xi\,\bigl(1-\tfrac{275}{4}\xi\bigr)\ \mathrm{eV}.
\end{equation}
In SI, the first order gives $3.7978\times10^{-23}$~J, i.e.\ $2.7507$~K and $+0.93\,\%$; the
second order gives $2.72549$~K and $+0.0002\,\%$ \KM. The agreement is real, its meaning is not:
three observations show that it depends on the unit.

\emph{First}, the relation holds only in eV. Read $\frac{16}{9}\xi$ in meV, in kelvin or in units
of $m_e$, and it is off by factors of $10^3$, $10^4$ or $5\times10^5$ \KM. Written as a ratio it
reads $T_{\mathrm{CMB}}/m_e=\frac{16}{9}\xi/510\,999$; the number $510\,999$ is $m_e$ in eV, i.e.\
the anchor itself, not an expression in $\xi$.

\emph{Second}, redefining $\alpha$ changes nothing. With $\alpha=1$ the charge unit grows by
$1/\sqrt\alpha=11.706$ and the potential shrinks by the same factor; energies stay the same \BM.
$T_{\mathrm{CMB}}$ contains no charge, so $\alpha$ can neither create nor remove the relation.

\emph{Third}, the eV is set, not measured \QM. Its size goes back to the practical volt of 1881
($10^8$ electromagnetic CGS units, chosen after the voltage of the Daniell cell), and thus to the
centimetre, gram and second. Every later stage -- the international volt of 1908, absolute units
in 1948, the Josephson constant in 1990 -- kept this size. Since 20~May 2019, with the fixed
elementary charge, exactly $1\ \mathrm{eV}=1.602176634\times10^{-19}$~J (SI Brochure, 9th ed.).
The CMB was discovered in 1965; it cannot have entered the size of the eV.

Then there is the coefficient. Exactly required would be $68.766$; $275/4=68.75$ is the nearest
fraction with denominator~4. Fractions $n/4$ lie $3.3\times10^{-5}$ apart in $T$, so a fitting
one can always be found to within $\pm1.7\times10^{-5}$. The second order lies $0.01\sigma$ from
the central value. Beyond the measurement precision an agreement proves nothing; a correct formula
would typically lie about $1\sigma$ off, and it comes this close only with a probability of about
$1\,\%$. This suggests that the coefficient was chosen at the central value. The classification of P11 in R137~(i) stands: unit-dependent agreement,
coefficient not derived \SM.

For the bookkeeping this means: a formula of the form ``quantity $=$ number times eV'' uses the
unit eV as an anchor. That is allowed as long as it is the only one. The mass ladder, however,
uses $m_e$; both together require that $m_e/\mathrm{eV}$ follow from $\xi$. The only proximity is
$\mathrm{eV}/m_e=1.271\,\xi^{3/2}$, $0.17\,\%$ from $(4/\pi)\,\xi^{3/2}$ \KM; that would derive
the size of the volt from $\xi$ and is not viable \XM.

% ============================================================
\section{Other reference energies}
% ============================================================

The basic form fixes an energy unit of its own: from $\xi E_0^2=1$ follows $u=\sqrt{\xi}\,E_0$,
so with $E_0=\sqrt{m_em_\mu}$, $u=84.85$~keV. In this unit
$T_{\mathrm{CMB}}=2.768\times10^{-9}\,u=0.0876\cdot\frac{16}{9}\xi^2\,u$ \KM, without a clean
prefactor. Referred to the energy scale of the atom,
$T_{\mathrm{CMB}}/(\alpha^2m_e)=0.0647\,\xi$ \KM, likewise without a clean prefactor. Neither
reference yields a relation.

% ============================================================
\section{The basic form: $T_{\mathrm{CMB}}/m_e$}
% ============================================================

In natural units $T_{\mathrm{CMB}}$ depends on the Hubble rate, the Planck energy and the share of
photons in the critical density. With $\rho_\gamma=\frac{\pi^2}{15}T^4$ and
$\rho_{\mathrm{crit}}=3H_0^2E_P^2/(8\pi)$,
\begin{equation}
\frac{T_{\mathrm{CMB}}}{m_e}=\Bigl(\frac{45\,\Omega_\gamma}{8\pi^3}\Bigr)^{1/4}
\Bigl(\frac{H_0E_P}{m_e^2}\Bigr)^{1/2}.
\end{equation}
With the corpus forms $H_0=\frac{\pi}{2}\xi^{10}m_e$ (Docs.~165, 279) and
$m_e/E_P=\frac{4/3}{5\pi}\xi^{11/2}$ (Doc.~383), $H_0E_P/m_e^2=\frac{15\pi^2}{8}\xi^{9/2}$ \BM. If
$\Omega_\gamma$ is proportional to $\xi$, then $T_{\mathrm{CMB}}/m_e\propto\xi^{5/2}$. Measured:
\begin{equation}
\frac{T_{\mathrm{CMB}}}{m_e}=K\,\xi^{5/2},\qquad K=2.23897 \quad\KM.
\end{equation}
This explains the exponent $2.41$ read off without a prefactor. The question is thereby reduced to
one number, equivalent to $\Omega_\gamma/\xi$.

Simple candidates lie outside the measurement precision: $\sqrt5$ ($-0.13\,\%$), $9/4$
($+0.49\,\%$), $64/(9\pi)$ ($+1.10\,\%$) \XM. Inside lies
\begin{equation}
\boxed{\frac{T_{\mathrm{CMB}}}{m_e}=(8\pi)^{1/4}\,\xi^{5/2}}
\qquad K=2.23903,\ +0.0025\,\%,\ 0.12\sigma \quad\KM.
\end{equation}
$8\pi$ is the coupling in $8\pi G$, which motivates the candidate but does not justify it \SM.

As a check against chance, all $3370$ expressions $(p/q\cdot\pi^k)^{1/n}$ with $p,q\le30$,
$|k|\le3$, $n\le4$ in the range $1.5$ to $3.5$ were searched. Three fall within FIRAS; $1.6$ are
expected by chance \KM. Restricted to $p,q\le10$ (383 expressions), only $(8\pi)^{1/4}$ remains,
with $0.2$ hits expected by chance. The candidate is thus the simplest, but not the only one; a
coincidence is not excluded.

% ============================================================
\section{Consequences of the candidate}
% ============================================================

\emph{Relation to $H_0$.} Equivalently $T_{\mathrm{CMB}}^4=8\pi\,\xi^{10}m_e^4$, with the same
exponent~10 as in $H_0$. Together with $H_0=\frac{\pi}{2}\xi^{10}m_e$, $\xi$ drops out:
\begin{equation}
\boxed{T_{\mathrm{CMB}}^4=16\,H_0\,m_e^3}\qquad\BM.
\end{equation}
From the measured $T_{\mathrm{CMB}}$ and $m_e$ follows $H_0=66.81\pm0.06$~km/s/Mpc \KM. This lies
$1.2\sigma$ from Planck ($67.4\pm0.5$) and far from SH0ES. The uncertainty comes from FIRAS alone
and is about eight times smaller than that of direct $H_0$ measurements.

\emph{Photon density.} With the T0 values of $H_0$ and $E_P$,
$\Omega_\gamma=\frac{4096}{10125}\,\xi=0.4045\,\xi$ \BM; with the measured $E_P$ instead of the T0
form, $0.4154\,\xi$.

\emph{$L_\xi$ without the CMB.} Inserting the candidate into the definition of $L_\xi$,
$T_{\mathrm{CMB}}$ drops out:
\begin{equation}
L_\xi=\Bigl(\frac{15}{8\pi^3}\Bigr)^{1/4}\xi^{-9/4}\,\bar\lambda_e=100.24\,\mu\mathrm{m}\qquad\KM.
\end{equation}
$L_\xi$ is thus a length built from $\xi$ and the Compton wavelength of the electron. The ratio
$\pi^2/(720\xi)=102.8$ between the Casimir and CMB densities at $L_\xi$ remains an identity; the
circularity, however, moves from $L_\xi$ to the prefactor $(8\pi)^{1/4}$.

% ============================================================
\section{The $H_0$ forms compared}
% ============================================================

The corpus contains two forms: $E_H=E_0\,\xi^{41/4}$ (Docs.~026, 064; exponent calibrated to the
measured value, P35) and $H_0=\frac{\pi}{2}\xi^{10}m_e$ (Docs.~165, 279). They differ by $1\,\%$,
since $E_0\,\xi^{1/4}=0.9904\cdot\frac{\pi}{2}m_e$ \KM. In the $H_0$ measurement this goes
unnoticed; $T_{\mathrm{CMB}}$, however, is measured more than thirty times more precisely and
separates the forms via $T^4=16\,H_0m_e^3$:

\begin{center}
\small
\begin{tabular}{@{}p{5.2cm}p{2.0cm}p{2.0cm}p{3.6cm}@{}}
\toprule
$H_0$ form & $H_0$ (km/s/Mpc) & $T_{\mathrm{CMB}}$ (K) & Deviation \\
\midrule
$\frac{\pi}{2}\xi^{10}m_e$ & $66.821$ & $2.72555$ & $+0.003\,\%$, $0.1\sigma$ \KM \\
$E_0\,\xi^{41/4}$, $E_0=7.398$~MeV & $66.179$ & $2.71898$ & $-0.24\,\%$, $11\sigma$ \XM \\
$E_0\,\xi^{41/4}$, $E_0=\sqrt{m_em_\mu}$ & $65.731$ & $2.71436$ & $-0.41\,\%$, $20\sigma$ \XM \\
backwards from $T_{\mathrm{CMB}}$ & $66.81\pm0.06$ & measured & -- \\
\bottomrule
\end{tabular}
\end{center}

If the candidate holds, the form with $41/4$ is excluded by the CMB, and the transition to the
exponent~10 in Doc.~165 is confirmed. This confirmation is not independent: the prefactor
$(8\pi)^{1/4}$ was found at the measured value, and the formula already contains $\xi^{10}$. What
is shown is that the same $\xi^{10}$ structure carries both $H_0$ and $T_{\mathrm{CMB}}$, with the
simple prefactors $\pi/2$ and $8\pi$.

% ============================================================
\section{Comparison}
% ============================================================

\begin{center}
\small
\begin{tabular}{@{}p{4.4cm}p{2.6cm}p{2.0cm}p{4.2cm}@{}}
\toprule
Route & Anchor & Agreement & Classification \\
\midrule
$\frac{16}{9}\xi$~eV & eV & $+0.93\,\%$ & unit-dependent \SM \\
$\frac{16}{9}\xi(1-\frac{275}{4}\xi)$~eV & eV & $2\times10^{-6}$ & unit-dependent, fitted \SM \\
$T/u$, $u=\sqrt\xi E_0$ & $m_e$ & no prefactor & no relation \\
$T/(\alpha^2m_e)$ & $m_e$ & no prefactor & no relation \\
$\sqrt5$, $9/4$, $64/(9\pi)$ & $m_e$ & $0.13$ to $1.1\,\%$ & outside FIRAS \XM \\
$(8\pi)^{1/4}\xi^{5/2}$ & $m_e$ & $0.12\sigma$ & candidate \SM \\
$\Lambda$CDM & $T_{\mathrm{CMB}}$ itself & -- & measured value as input \QM \\
\bottomrule
\end{tabular}
\end{center}

% ============================================================
\section{What remains open}
% ============================================================

The circularity of $L_\xi$ is resolved only conditionally: it has been moved to one number. Open
are
\begin{itemize}
\item the prefactor $(8\pi)^{1/4}$, equivalent to $\Omega_\gamma=\frac{4096}{10125}\xi$:
motivated by $8\pi G$, not derived; with $1.6$ hits expected in the wide search space, a
coincidence is not excluded \SM;
\item the exponent~10 and the prefactor $\pi/2$ of $H_0$ (P20); the $T$ relation inherits them,
since $5/2=10/4$ \SM;
\item why the photon density today is a fixed fraction of $\xi$; in an expanding model this would
be a coincidence of epoch, in a static one a statement that needs a derivation \SM;
\item the eV relation P11: unit-dependent, coefficient $275/4$ fitted; it remains as a historical
agreement but contributes nothing to the basic form (R137, R70) \SM.
\end{itemize}
The candidate can be tested through $H_0$: it requires $H_0=66.81\pm0.06$~km/s/Mpc. A
model-independent $H_0$ measurement with a precision of a few tenths of a per cent would confirm
or exclude it.

% ============================================================
\section{Conclusion}
% ============================================================

The eV relation agrees because the eV brings the electron mass along as a second anchor; in every
unit that FFGFT fixes itself, the agreement disappears. In the basic form,
$T_{\mathrm{CMB}}/m_e$ depends on a single number, and the simplest candidate,
$(8\pi)^{1/4}\xi^{5/2}$, meets FIRAS to $0.12\sigma$. It links $T_{\mathrm{CMB}}$ and $H_0$ as
$T^4=16\,H_0m_e^3$, fixes $L_\xi$ without the CMB, and excludes the old $H_0$ form with $41/4$.
It is not derived; the prefactor and, as for $H_0$, the exponent~10 remain open.

\vspace{1em}
\noindent\textit{Verification script:} \texttt{2/python/Dok388\_Skripte/pruef\_388\_cmb\_grundform.py}
--- 27/27.
